\section{Related Work}
\label{sec:related}

\paragraph{VLM Hallucination.}
Hallucination in large vision-language models---generating content inconsistent with
the visual input---has been extensively documented~\cite{li2023evaluating,guan2024hallusionbench}.
Existing work identifies hallucinations in object existence, attributes, spatial
relations, and counts. Most prior work treats hallucination as a generation artifact
and addresses it via training-time regularization or decoding
modifications~\cite{leng2024mitigating,huang2024opera}.
\bci{} instead reframes the problem as an \emph{epistemic} failure amenable to
inference-time structural correction, without changing model parameters.

\paragraph{Post-Hoc Verification and Correction.}
Woodpecker~\cite{yin2023woodpecker} is the closest prior work: it detects and corrects
hallucinations post-hoc by querying specialized models and retrieving image evidence.
\bci{} differs in two key respects: (1)~it corrects beliefs \emph{before} the final
answer is generated, closing the reasoning loop at the premise level; and (2)~it
quantifies the causal contribution of false premises via controlled ablation.
Chain-of-Verification (\textsc{CoVe})~\cite{dhuliawala2024chain} generates targeted
re-check questions after an initial response and revises accordingly; this is a
complementary self-refinement approach, but it does not decompose failures into
premise \vs reasoning types. Our matched in-house evaluation of a CoVe-lite baseline
confirms that \bci{} outperforms this strategy on GQA.

\paragraph{Inference-Time Intervention.}
A growing literature explores inference-time methods for improving LLM and VLM
reliability. Self-consistency~\cite{wang2023self} aggregates multiple samples;
resampling-based methods~\cite{cobbe2021training} filter by answer frequency. Visual
Contrastive Decoding~\cite{leng2024mitigating} modifies logit distributions to
suppress hallucination. These approaches operate at the \emph{output distribution}
level, whereas \bci{} intervenes at the \emph{intermediate belief} level, offering
interpretability and fine-grained error attribution.

\paragraph{Scene Graph Grounding.}
GQA~\cite{hudson2019gqa} provides structured scene-graph annotations that make formal
claim verification tractable. We use these as a deterministic verifier in Phase~1
experiments, deliberately separating verifier design from the core intervention
hypothesis. Extending \bci{} to neural or learned verifiers is a natural next step.

\paragraph{Error Decomposition.}
Prior work on VQA error analysis~\cite{antol2015vqa,goyal2017making} typically operates
at the answer level. Our premise/reasoning decomposition provides a finer-grained
taxonomy more suitable for targeted intervention design.
